%%%%%%%%%%%%%%%%%%%%%%% CHAPTER - 9 %%%%%%%%%%%%%%%%%%%%
\chapter{Conclusion}
\label{C9} %%%%%%%%%%%%%%%%%%%%%%%%%%%%
%\graphicspath{{Figures/Chapter-6figs/PDF/}{Figures/Chapter-6figs/}}
% \noindent\rule{\linewidth}{2pt}
%%%%%%%%%%%%%%%%%%%%%%%%%%%%%%%%%%%%%%%%%%%%%%%%%%%%%%%%%%%%%%%%%%%%%%%%%%%%%%%%%%
The AI Cricket Coach system successfully demonstrates the practical integration of artificial intelligence, computer vision, and mobile application development to create an intelligent cricket training assistant. The project was designed to detect and classify cricket batting shots in real time using the YOLOv8 object detection model and to evaluate shot performance using an XGBoost regression model. By combining deep learning–based pose and object detection with machine learning–based performance scoring, the system provides both shot recognition and technical assessment within a unified platform.The implementation of a client–server architecture ensures efficient distribution of computational tasks, where the Android application handles image capture and user interaction, while the Flask backend server processes detection and scoring operations. The use of a well-structured dataset with proper annotation, augmentation, and validation contributed significantly to achieving reliable detection accuracy. The inclusion of multiple detection outcomes such as Straight Drive, Cut Shot, Pull Shot, Sweep Shot, Unknown stance, and No Person Detected improves system robustness and prevents incorrect classifications. Additionally, the integration of real-time feedback and Text-to-Speech functionality enhances user engagement and makes the application practical for cricket practice sessions.Through extensive testing and validation, the system demonstrated accurate shot detection, consistent performance scoring, and stable client–server communication. The modular and scalable architecture allows for future enhancements such as live video analysis, cloud deployment, multi-user performance tracking, and advanced biomechanical analytics.
In conclusion, the AI Cricket Coach project proves that artificial intelligence can effectively assist in sports training by providing data-driven insights, automated technical analysis, and instant feedback. The system has the potential to support players, coaches, and training institutions by making cricket coaching more accessible, objective, and technology-driven.